(1967) (for summary, see Wheeler 1966), by Misner and Zapolsky (1967), and by Chau
(1967) suggest that the non-radial pulsations of a neutron star should damp out due
to the emission of gravitational waves in a time $\tau \lesssim 1$ sec, and that, consequently, the
gravitational waves emitted by a neutron star that is formed in a supernova explosion in
our own Galaxy might be detectable by the apparatus of Joseph Weber (see Weber
1964, 1967). In order to make these rough estimates of the gravitational-radiation flux
from a neutron star more precise, one needs the detailed theory of non-radial pulsations
of relativistic stellar models.

The theory of non-radial pulsations is important not only because of its astrophysical
implications but also because of its key status within the framework of the theory of
gravitational radiation. We know of no self-consistent analysis to date, within the
framework of the full theory of relativity, of the emission of gravitational waves by a
dynamical system---and of the consequent damping of that system's motion. When completed,
numerical computations based upon the equations derived here will provide
such an analysis.

The presentation of the theory of non-radial pulsations is divided into six sections.
In \S~II perturbations with arbitrary time dependences are analyzed into tensorial
spherical harmonics, and the equations of motion and of radiation, governing evolution of
the perturbations, are derived and discussed. Pulsation is absent in perturbations of
``odd parity,'' so the subsequent analysis is confined to ``even-parity'' perturbations. In
\S~III the even-parity perturbations are analyzed into complex normal modes with various
mixtures of incoming and outgoing gravitational radiation. The eigenproblems---including
differential equations and boundary conditions---is formulated for the complex
normal modes with spherical harmonic index $l \geq 2$; and the form of the gravitational
waves far from the star is discussed. Section~IV contains a discussion of the relationship
between the complex normal modes and the real pulsations of a relativistic stellar model.
Real, ``quasi-normal'' pulsations which are initiated at a particular moment of time---
and only pump outgoing gravitational waves with a sharp wave front---are expanded in
terms of complex normal modes. The frequencies, eigenfunctions, and damping times of
the complex normal solutions are related to properties of the complex normal mode with
purely outgoing radiation. Section~V is a description of methods for calculating numerically
the frequencies and eigenfunctions of the complex normal modes with purely
outgoing radiation; and \S~VI is a discussion of the directions in which this analysis
should be pushed further. In order to make the paper readable, we have confined to
appendices the derivations of most of the equations.

Some of the results presented in this paper have been derived independently but not
published by S.\ Chandrasekhar and by H.\ Zapolsky.

\section{ARBITRARY, SMALL PERTURBATIONS}

\subsection*{(a) The Equilibrium Configuration}

Throughout this paper we shall use the conventions and notation of Thorne (1966,
1967), ``Relativistic Stellar Structure and Dynamics''; cited henceforth as RSSD) except
that we shall adopt geometrized units ($c = G = 1$). We shall omit the
asterisks ($*$) used in RSSD for geometrized quantities, and we shall use the symbol $\nu$ in
place of $2\Phi$ for $\ln g_{tt}$.

As is discussed in RSSD, a relativistic equilibrium configuration can be described by a
coordinate system $(t, r, \theta, \phi)$, with respect to which the geometry of spacetime is given
by the line element
\begin{equation}
  ds^{2} = -(e^{\nu})\,(dt)^{2} + e^{\lambda}\,(dr)^{2} + r^{2}(d\theta^{2} + \sin^{2}\theta\,d\phi^{2})\;.
\label{eq:1}
\end{equation}

Here $\nu$ and $\lambda$ are functions of the radius $r$, and $m(r)$ is defined as the ``interior
mass'' $m(r)$, by
\begin{equation}
  e^{-\lambda} = 1 - 2m/r\;.
\label{eq:2}
\end{equation}
